\section{从 $m=6$ 推进到跨尺度折叠：保持核不变的分辨率控制}\label{sec:multiscale}

\subsection{一般 $\Fold_m$ 的迭代性（提纲）}
对一般窗口长度 $m$，仍有
\[
  \abs{\Omega_m}=2^m,\qquad \abs{X_m}=F_{m+2},
\]
并可定义 $\Fold_m$ 为 Zeckendorf 位串截断至前 $m$ 位得到 $X_m$ 元素。其预像结构可用“窗口外 Fibonacci 权重的可容许子集”刻画，并给出增长界与平均退化
\[
  \frac{2^m}{F_{m+2}}\approx \left(\frac{2}{\varphi}\right)^m.
\]

\subsection{保持 $m=6$ 核不变的跨尺度方案（接口层描述）}
我们给出一个可操作的跨尺度方案：把更大尺度的信息通过控制指令 $u$ 与路径记忆（在 $\widetilde G_6$ 上的 uplift/shift 轨迹）来表达，而稳定型字母表保持为 $X_6$ 的 $21$ 个类型。形式化地，这相当于用 $\widetilde G_6$ 的路径空间（或其有限状态扩张）承载更大 $m$ 的有效状态；而 $m$ 的“提升/降低”变为图上的控制策略，而非更换核本身。

